\documentclass[11pt,a4paper]{article}
% --- PAKETE FÜR MATHE UND TYPOGRAFIE ---
\usepackage[utf8]{inputenc}
\usepackage[T1]{fontenc}
\usepackage{lmodern}
\usepackage[ngerman]{babel}
\usepackage{amsmath,amssymb,amsthm,amsfonts}
\usepackage{mathtools}
\usepackage{bm}
\usepackage{booktabs}
\usepackage{array}
\usepackage{graphicx}
\usepackage{microtype}
\usepackage{hyperref}
\usepackage{geometry}
\usepackage{verbatim}
\geometry{
	a4paper,
	left=25mm,
	right=25mm,
	top=30mm,
	bottom=30mm
}
\hypersetup{
	colorlinks=true,
	linkcolor=blue,
	citecolor=blue,
	urlcolor=blue
}
% --- THEOREM-UMGEBUNGEN ---
\newtheorem{theorem}{Theorem}[section]
\newtheorem{lemma}[theorem]{Lemma}
\newtheorem{proposition}[theorem]{Proposition}
\theoremstyle{definition}
\newtheorem{definition}[theorem]{Definition}
\newtheorem{remark}[theorem]{Bemerkung}
% --- BENUTZERDEFINIERTE MATHEMATISCHE MAKROS ---
\newcommand{\Oalg}{\mathbb{O}}
\newcommand{\ImO}{\operatorname{Im}(\mathbb{O})}
\newcommand{\Assoc}{\operatorname{Assoc}}
\newcommand{\Jac}{\operatorname{J}}
\newcommand{\RePart}{\operatorname{Re}}
\newcommand{\ImPart}{\operatorname{Im}}
% --- TITEL UND AUTOREN-METADATEN ---
\title{\textbf{Eine Vier-Ebenen-Systemarchitektur zur formalen und numerischen Trennung oktonionischer Kompositionsalgebren von analytischen Zeta-Schichten}}
\author{%
	\textbf{Thomas Hoffbauer}\thanks{Computer Institut Bamberg (CIB), Deutschland. E-Mail: \texttt{research@cib-bamberg.de}} \\[2mm]
	\small\textit{Projekt \#Energiedoku -- Preprint-Reihe: Formalisierung und Renormierung (eabc-renorm)}
}
\date{3. August 2026}
\begin{document}
	\maketitle
	\begin{abstract}
		Wir präsentieren eine typologisch und epistemisch rigorose Vier-Ebenen-Architektur zur Einbettung der Oktonionenalgebra $\Oalg$ in gittertheoretische und komplex-analytische Kontexte. Ausgehend von der expliziten Vektorraumzerlegung des Oktonionenassoziators $\Assoc(x,y,z) = (xy)z - x(yz)$ auf dem sieben-dimensionalen Imaginärraum $\ImO \cong \mathbb{R}^7$ zeigen wir die exakte Proportionalität des Kreuzprodukt-Jacobiators $J_\times = \frac{3}{2}\Assoc$ sowie die Unterscheidung zwischen der multiplikativen Norm $N(x) = \|x\|^2$ und der gittertheoretischen quadratischen Form $Q(x) = \frac{1}{2}\|x\|^2$. Wir beschreiben die Formalisierung dieser Trennung im Lean-4-Kernel unter Vermeidung assoziativer Typklassen (\texttt{[Ring A]}) sowie die reproduzierbare numerische Verifikation mittels einer tautologie-freien Python-Audit-Suite. Schließlich belegen wir anhand eines quaternionischen Gegenbeispiels, dass der alternierte ternäre Kommutator (Ternutator) keine universelle Proportionalitätsidentität $T = c \Assoc$ erfüllt, und verankern strikte Non-Claim-Guardrails gegen kategoriale Fehlschlüsse in der mathematischen Physik.
	\end{abstract}
	\noindent\textbf{Keywords:} Oktonionen, Malcev-Algebra, $G_2$-Symmetrie, Assoziator, Epstein-Zeta-Funktion, Lean 4 Formalisierung, Formal Governance.
	\section{Einleitung}
	Die Untersuchung nicht-assoziativer Divisionsalgebren, insbesondere der Reellen Oktonionen $\Oalg$, nimmt in der modernen mathematischen Physik und algebraischen Topologie eine zentrale Rolle ein \cite{baez2002octonions}. Bei der Konstruktion von Modellen, die oktonionische Geometrie mit spektralen Eigenschaften diskreter Gitter (wie dem $E_8$-Gitter) oder komplex-analytischen Zählfunktionen (Epstein-Zeta-Funktionen) verknüpfen, entstehen häufig kategoriale Überdehnungen. Nicht-assoziative Identitäten werden dabei fälschlicherweise unmittelbar mit physikalischen Eichfeldern, Nambu-Dynamiken oder Paritätsverletzungen identifiziert.
	Das Ziel dieser Arbeit ist die Etablierung einer unumstößlichen
	\textbf{Vier-Ebenen\-System\-architektur} im Rahmen des Projekts
	\texttt{eabc-renorm} (Pull Request \#143, Commit \texttt{7ff7c62}). Diese Architektur trennt die lineare Gitterstruktur, die $G_2$-Malcev-Geometrie, die alternative Kompositionsalgebra sowie die komplex-analytische Zeta-Schicht formal-logisch (Lean 4) und numerisch (Python/SageMath).
	\section{Die Vier-Ebenen-Systemarchitektur}
	Um unbeabsichtigte Implikationen zu vermeiden, unterteilen wir den mathematischen Formalismus streng in vier epistemische Ebenen (siehe Tabelle~\ref{tab:four_levels}).
	\begin{table}[h!]
		\centering
		\caption{Systemarchitektur der vier epistemischen Ebenen mit zugehörigen Symmetrien und Formalisierungsklassen.}
		\label{tab:four_levels}
		\vspace{2mm}
		\resizebox{\textwidth}{!}{%
			\begin{tabular}{@{}lllll@{}}
				\toprule
				\textbf{Ebene} & \textbf{Algebraische Struktur} & \textbf{Kern-Operation} & \textbf{Symmetrie} & \textbf{Status-Typ} \\
				\midrule
				\textbf{I. Linear} & $(\mathbb{R}^8, +, \text{Einbettung})$ & Vektoraddition $x+y$ & Diskrete Automorphismen & \texttt{[I]} Rank-8-Interface \\
				\textbf{II. Malcev} & $V = \ImO \cong \mathbb{R}^7$ & Kreuzprodukt $x \times y$ & $G_2 \subset SO(7)$ & \texttt{[I]} Normed Cross / \texttt{[A-ext]} \\
				\textbf{III. Alternativ} & Unitale Algebra $\Oalg$ & Vollmultiplikation $x \cdot y$ & $G_2$ (Algebra-Aut.) & \texttt{[I]} Alternative Algebra \\
				\textbf{IV. Analytisch} & Zeta-Schicht $Z_L(s)$ & Form $\omega_Z = \frac{Z'}{Z}ds$ & Komplex-Analytisch & \texttt{[I]} Analytische Schicht \\
				\bottomrule
			\end{tabular}%
		}
	\end{table}
	\subsection{Ebene I: Linear- und Gitterstruktur}
	Ebene I stellt das Trägermedium bereit. Ein achtdimensionaler reeller Vektorraum $V \cong \mathbb{R}^8$ enthält eine injektive Einbettung $\iota \colon \mathbb{Z}^8 \hookrightarrow V$. Wir betonen, dass die reine Linearstruktur $V$ noch kein $E_8$-Gitter definiert; hierzu bedarf es der Angabe der positiv definiten, geraden, unimodularen Gram-Matrix.
	\subsection{Ebene II: \texorpdfstring{$G_2$}{G2}-Kreuzprodukt und Malcev-Struktur}
	Auf dem Raum der rein imaginären Oktonionen $\ImO \cong \mathbb{R}^7$ wird das bilinear alternierende Kreuzprodukt über den Kommutator definiert:
	\begin{equation}
		x \times y \coloneqq \frac{1}{2} \ImPart(x \cdot y - y \cdot x) = \frac{1}{2} [x, y]_{\mathrm{comm}}.
	\end{equation}
	Die Kommutatoralgebra der Oktonionen ist keine Lie-Algebra, sondern bildet bezüglich dieses Produkts eine einfache Malcev-Algebra.
	\subsection{Ebene III: Alternative Kompositionsalgebra}
	Die volle Oktonionenmultiplikation vereint den symmetrischen Skalaranteil und den antisymmetrischen Vektoranteil:
	\begin{equation}
		x \cdot y = -\langle x, y \rangle + x \times y \quad \forall\, x,y \in \ImO.
	\end{equation}
	$\Oalg$ bildet eine normierte alternative Algebra, deren Nichtnull-Elemente $\Oalg^\times$ einen Moufang-Loop bilden.
	\subsection{Ebene IV: Analytische Zeta-Schicht}
	Aus der gittertheoretischen quadratischen Form $Q(x)$ wird die Epstein-Zeta-Funktion $Z_L(s) = \sum_{x \in L \setminus \{0\}} Q(x)^{-s}$ konstruiert. Die logarithmische Differentialform $\omega_Z = \frac{Z_L'(s)}{Z_L(s)} \, ds$ fungiert als analytische Zählform.
	\begin{proposition}[Entkopplung der Zeta-Linie]
		Die meromorphe Fortsetzung von $Z_L(s)$ ist eine Eigenschaft der diskreten Gitterstruktur $(L, Q)$ und folgt \textbf{nicht} aus der Oktonionenmultiplikation auf Ebene III:
		\begin{equation}
			\begin{array}{ccccc}
				(L, Q) &\longrightarrow& \Theta_L, \; Z_L &\longrightarrow& \dfrac{Z_L'}{Z_L} \, ds \\[2mm]
				\mathbb{O} &\longrightarrow& \operatorname{Im}\mathbb{O}, \; \times, \; \operatorname{Assoc} &\longrightarrow& \text{Malcev-$G_2$-Geometrie}.
			\end{array}
		\end{equation}
	\end{proposition}
	\section{Algebraische Identitäten des Assoziators}
	\begin{definition}[Oktonionenassoziator]
		Für $x, y, z \in \Oalg$ ist der Assoziator definiert als
		\begin{equation}
			\Assoc(x,y,z) \coloneqq (xy)z - x(yz).
		\end{equation}
		Infolge der Alternativität von $\Oalg$ ist $\Assoc$ eine total schiefsymmetrische, alternierende 3-Operation.
	\end{definition}
	\begin{theorem}[Exakte Vektorraumzerlegung]
		Für rein imaginäre Oktonionen $x, y, z \in \ImO$ zerfällt der Assoziator exakt in Skalarprodukt- und Kreuzproduktkomponenten:
		\begin{equation}
			\label{eq:assoc_decomp}
			\Assoc(x,y,z) = -\langle x,y \rangle z + \langle y,z \rangle x + (x \times y) \times z - x \times (y \times z).
		\end{equation}
	\end{theorem}
	\begin{proof}
		Unter Nutzung von $xy = -\langle x,y \rangle + x \times y$ berechnet man:
		\begin{align}
			(xy)z &= -\langle x \times y, z \rangle - \langle x,y \rangle z + (x \times y) \times z, \\
			x(yz) &= -\langle x, y \times z \rangle - \langle y,z \rangle x + x \times (y \times z).
		\end{align}
		Wegen der $G_2$-Invarianz der skalaren 3-Form $\varphi(x,y,z) = \langle x \times y, z \rangle = \langle x, y \times z \rangle$ heben sich die skalaren Dreifachprodukte in der Differenz $\Assoc(x,y,z) = (xy)z - x(yz)$ identisch auf. Es verbleibt Gl.~\eqref{eq:assoc_decomp}.
	\end{proof}
	\begin{remark}[Alternativität via Kreuzprodukt-Identität]
		Für $y = x$ reduziert sich Gl.~\eqref{eq:assoc_decomp} auf $\Assoc(x,x,z) = -\|x\|^2 z + \langle x,z \rangle x - x \times (x \times z)$. Durch Anwendung der 7D-Kreuzprodukt-Identität
		\begin{equation}
			x \times (x \times z) = \langle x,z \rangle x - \|x\|^2 z
		\end{equation}
		folgt $\Assoc(x,x,z) = 0$. Die Alternativität verschwindet im Imaginärraum somit nicht trivial, sondern durch ein exaktes algebraisches Aufheben.
	\end{remark}
	\begin{theorem}[Kreuzprodukt-Jacobiator]
		Der Jacobiator des $G_2$-Kreuzprodukts $J_\times(x,y,z) \coloneqq (x \times y) \times z + (y \times z) \times x + (z \times x) \times y$ erfüllt:
		\begin{equation}
			J_\times(x,y,z) = \frac{3}{2} \Assoc(x,y,z) \quad \forall\, x,y,z \in \ImO.
		\end{equation}
	\end{theorem}
	\begin{proposition}[Norm-Distinktion]
		Zur Vermeidung von Skalierungsfehlern unterscheiden wir streng zwischen der multiplikativen Norm $N(x) \coloneqq \|x\|^2$ und der gittertheoretischen quadratischen Form $Q(x) \coloneqq \frac{1}{2}\|x\|^2$. Es gilt:
		\begin{equation}
			N(xy) = N(x) N(y), \qquad Q(xy) = 2 \, Q(x) Q(y).
		\end{equation}
	\end{proposition}
	\section{Lean 4 Formalisierung und Typologie}
	Die Implementierung im Theorem Prover Lean 4 (Modul \texttt{FourLevelArchitecture.lean}) stellt sicher, dass abstrakte Schnittstellen nicht fälschlich als bewiesene Kernel-Sätze ausgegeben werden.
	\subsection{Vermeidung der Typen-Falle \texttt{[Ring A]}}
	Ein Lean-\texttt{Ring} beinhaltet per Konstruktion das Assoziativgesetz $(xy)z = x(yz)$. Eine Formalisierung von $\Oalg$ unter \texttt{[Ring A]} würde trivialerweise $\Assoc(x,y,z) = 0$ erzwingen. Ebene III wurde daher als \texttt{Level3\_UnitalAlternativeAlgebra} ohne \texttt{[Ring A]} modelliert:
	\begin{verbatim}
		structure Level3_UnitalAlternativeAlgebra (A : Type*)
		[AddCommGroup A] [Module R A] where
		one : A
		mul : A -> A -> A
		left_alternative  : forall x y, mul (mul x x) y = mul x (mul x y)